% !TeX root = ../navier-stokes-manuscript.tex
\clearpage
\section{One Terminal Datum and the Same-History Restart}\label{sec:TheMissingEstimateAndGlobalSmoothness}

The argument has arrived at the face where a smooth history could disappear.
On every strict subinterval the same fluid carries every finite derivative.
That interior smoothness is not the issue.  The issue is whether a fixed finite
rectangle keeps one bound while its cutoff moves all the way to \(T_*\).  If
the bound rises without limit, no terminal datum has been reached.  If the
whole-terminal condition in
\Cref{ass:BodyWholeTerminalRectangleCondition} holds, the adjacent rows reach
the face together, their overlaps identify one endpoint, and the local
equation can be restarted from the state at which the old history arrived.

This section follows that implication to its end.  It does not turn the
condition into an arbitrary-data estimate: the unresolved direction remains
\((u_0,\nu)\to D_*<\infty\), already isolated in
\eqref{eq:BodyFirstRectangleCriticalDissipationEquivalence}.

\subsection{The intersection reaches one terminal datum}\label{sub:NoFiniteSingularTime}

\begin{theorem}[Conditional whole-space continuation]\label{thm:MainGlobalRegularity}
Let \(\nu>0\), let \(u_0\in\Ssigma\), and let \((u,p)\) be the maximal
classical history of \eqref{eq:NavierStokesSystem}.  Assume its finite branch,
if one exists, satisfies
\Cref{ass:BodyWholeTerminalRectangleCondition}.  Then
\begin{equation}\label{eq:BodyGlobalSmoothnessConclusion}
  T_*=\infty,
  \qquad
  u,p\in C^\infty\!\bigl(\R^3\times[0,\infty)\bigr).
\end{equation}
For every \(t\geq0\),
\begin{equation}\label{eq:GlobalEnergyConclusion}
  \frac12\norm{u(t)}_2^2
  +\nu\int_0^t\norm{\nabla u(s)}_2^2\,ds
  =\frac12\norm{u_0}_2^2.
\end{equation}
\end{theorem}

\begin{proof}
Suppose instead that the maximal history ends at some finite \(T_*\).  The
assumed whole-terminal condition carries every finite pressure-complete
rectangle across \(I_*\).  Whenever two rectangles overlap,
\eqref{eq:BodyRectangleRestrictionCompatibility} says that they contain the
same derivatives of the same history.  The mixed intersection in
\Cref{prop:BodyCompatibleProjectiveAssembly} gives
\[
  u\in\bigcap_{r,m<\infty}H^r(I_*;H^m(\R^3)).
\]

Hold one spatial height \(m\) fixed.  The temporal row \(r=1\) places \(u\)
in \(H^1(I_*;H^m)\), so the actual history has an \(H^m\) limit as
\(t\uparrow T_*\).  A higher row cannot arrive at a different field: after it
is embedded into \(H^m\), both traces are limits of the same \(u(t)\).  Every
finite height reaches one divergence-free endpoint
\begin{equation}\label{eq:BodyCommonSmoothEndpoint}
  u_*:=\lim_{t\uparrow T_*}u(t)
  \in\Hinfty_\sigma(\R^3)
  \subset\bigcap_{m<\infty}H^m(\R^3).
\end{equation}
This is the endpoint furnished by \Cref{lem:CompleteSpatialRowEndpoint}.

The endpoint must remember how the history arrives, not only the velocity at
which it arrives.  Given
finite \(N\) and \(m\), the row \(r=N+1\) places
\(V_N=\partial_t^Nu\) in \(H^1(I_*;H^m)\).  It has a trace \(V_N^*\), and
spatial compatibility puts that trace in \(\Hinfty\).  The normalized pressure
recurrence and the velocity recurrence survive the limit:
\begin{align*}
  Q_N^*
  &=R_iR_j\sum_{a=0}^{N}\binom{N}{a}V_{a,i}^*V_{N-a,j}^*,
  \\
  V_{N+1}^*
  &=\nu\Delta V_N^*
  -\sum_{a=0}^{N}\binom{N}{a}(V_a^*\cdot\nabla)V_{N-a}^*
  -\nabla Q_N^*.
\end{align*}
Thus the entire left-hand velocity jet and the Riesz-normalized pressure jet
are generated by \(u_*\), as \Cref{prop:PressureCompleteEndpointJet} records.

Release \(u_*\) as initial data at time \(T_*\).  The local construction in
\Cref{lem:SmoothLocalRestart} gives a right-hand history on
\([T_*,T_*+\delta]\).  It begins from the same velocity and the same pressure
normalization, so its jet agrees with every left-hand trace.  The two sides are
one smooth solution.  Strong uniqueness identifies the right-hand piece with
the continuation of the history generated by \(u_0\), contradicting the
maximality of \(T_*\).

The assumed finite branch is impossible, so \(T_*=\infty\).  The normalized
pair is smooth on every finite time interval.  The kinetic-energy identity in
\Cref{prop:ClassicalKineticEnergyIdentity} now holds on every initial segment
of this global history and gives \eqref{eq:GlobalEnergyConclusion}.  In
particular,
\[
  \sup_{t\geq0}\int_{\R^3}|u(x,t)|^2\,dx
  \leq
  \int_{\R^3}|u_0(x)|^2\,dx
  <\infty.
\]
The rapid decay of \(u_0\) pays the last finiteness.  Under the whole-terminal
condition, this is the global history requested in
\eqref{eq:WholeSpaceRegularityTarget}.
\end{proof}

The full intersection carries the complete jet to the restart.  Its lowest
rectangle already contains the shorter endpoint contradiction that explains
why the first missing payment is decisive.

\subsection{What the first rectangle already knows}\label{sub:FromFiniteRectanglesToClassicalEstimates}

Return to the first rectangle, \((N,M)=(0,1)\).  Under the same terminal
condition, its adjacent rows are
\[
  u\in L^2(I_*;H^2),
  \qquad
  u_t\in L^2(I_*;L^2).
\]
Their trace reaches the space between them:
\[
  u\in C([0,T_*];H^1).
\]
Since \(H^1(\R^3)\hookrightarrow L^3(\R^3)\), this gives
\begin{equation}\label{eq:EndpointCriterionFromLowRectangle}
  \sup_{0\leq t\leq T_*}\norm{u(t)}_3<\infty.
\end{equation}
The endpoint theorem of Escauriaza, Seregin, and \v{S}ver\'ak says that a
finite maximal time forces the opposite behaviour
\parencite{EscauriazaSereginSverak2003}:
\begin{equation}\label{eq:FiniteTimeLThreeEscape}
  T_*<\infty
  \quad\Longrightarrow\quad
  \sup_{0<t<T_*}\norm{u(t)}_3=\infty.
\end{equation}
The rectangle and the endpoint theorem concern the same velocity, the same
norm, and the same terminal interval.  Their conclusions cannot coexist.  The
lowest rectangle is already enough to rule out the alleged last time; the
higher rectangles do different work by carrying the full jet into the smooth
restart.

We can now return to the physical history without beginning another proof.  The
first rectangle has reached the critical velocity norm.  The next rectangle
reaches the motion of the critical height and the gradient that stretches the
vortex.  These consequences tell us what the terminal condition has done to
the narrowing mechanism with which the paper began.

\subsection{The critical-height roof}\label{sub:ClassicalContinuationCriteriaInsideFiniteRectangles}

Once the \(H^1\) trace reaches the terminal face, the critical height cannot keep
climbing underneath it.  Indeed,
\(H=\tfrac12\norm{\Lambda^{1/2}u}_2^2\), so
\begin{equation}\label{eq:BodyCriticalHeightRoof}
  H_{\max}
  :=\sup_{0\leq t\leq T_*}H(t)
  \leq\frac12\sup_{0\leq t\leq T_*}\norm{u(t)}_{H^1}^2
  <\infty.
\end{equation}
Start at any positive height \(H_0\), double it successively by
\(L_k:=2^kH_0\), and count the levels the same history actually reaches:
\[
  N_{\mathrm{levels}}
  :=
  \#\bigl\{k\in\Nzero:\text{there is a }t\in[0,T_*]\text{ with }H(t)=L_k\bigr\}.
\]
The roof leaves room for only finitely many such doublings:
\begin{equation}\label{eq:BodyFiniteZenoOctaveCount}
  N_{\mathrm{levels}}
  \leq
  \begin{cases}
    0,&H_{\max}<H_0,\\
    1+\left\lfloor\log_2(H_{\max}/H_0)\right\rfloor,
      &H_{\max}\geq H_0.
  \end{cases}
\end{equation}
This conclusion never asks how quickly one level is crossed.
The summable heat clocks from Section~\ref{sec:CriticalHeightAndDerivativeTower} cannot
assemble an infinite ladder because the height itself runs out of levels.
A width might still shrink while the height stays bounded, so the geometric
picture has one more question to answer.

The next rectangle, \(\cR_{0,2}\), follows the motion beneath this roof rather
than recording only its maximum:
\begin{equation}\label{eq:LowRectangleZeroTwo}
  u\in L^2(I_*;H^3),
  \qquad
  u_t\in L^2(I_*;H^1),
  \qquad
  u\in C([0,T_*];H^2).
\end{equation}
The two adjacent rows now contain both factors in the derivative
\[
  H'(t)
  =\operatorname{Re}
  \pair{\Lambda^{1/2}u(t)}{\Lambda^{1/2}u_t(t)}_{L^2},
\]
and Cauchy--Schwarz pays its total variation:
\begin{equation}\label{eq:CriticalHeightTotalVariationBound}
  \int_0^{T_*}|H'(t)|\,dt
  \leq
  \norm{\Lambda^{1/2}u}_{L^2_tL^2_x}
  \norm{\Lambda^{1/2}u_t}_{L^2_tL^2_x}
  <\infty.
\end{equation}
The same rectangle has \(E_{3/2}\in L^1(I_*)\).  The completed critical row
\[
  \mathcal P_{1/2}=H'+2\nu E_{3/2}
\]
now places the signed production itself in \(L^1(I_*)\), and integration
records its accumulated action:
\begin{equation}\label{eq:BodySignedCriticalPrefix}
  \int_0^t\mathcal P_{1/2}(s)\,ds
  =H(t)-H(0)+2\nu\int_0^tE_{3/2}(s)\,ds.
\end{equation}
The two terms on the right are finite, so
\begin{equation}\label{eq:BodySignedCriticalPrefixBound}
  A_*(T_*)
  :=\sup_{0<t<T_*}\int_0^t\mathcal P_{1/2}(s)\,ds
  <\infty.
\end{equation}
The upper prefix in \eqref{eq:BodySignedCriticalPrefixBound} is already
equivalent to the lowest rectangle by
\eqref{eq:BodyFirstRectangleSignedPrefixEquivalence}; it is the same interface,
not a free upstream estimate.  The stronger \(\cR_{0,2}\) control adds finite
total variation of \(H\) and absolute time integrability of the completed
critical row.  Sharp rises and falls remain possible, but neither their total
height variation nor their signed accumulation can diverge before \(T_*\).

For the vortex, the decisive entry in \(\cR_{0,2}\) is
\(u\in L^2(I_*;H^3)\).  In three dimensions
\(H^3\hookrightarrow W^{1,\infty}\), so
\[
  \int_0^{T_*}\norm{\nabla u(t)}_\infty\,dt
  \leq
  T_*^{1/2}\norm{u}_{L^2(I_*;H^3)}<\infty.
\]
Both vorticity and strain are parts of this gradient.  Their accumulated action
before \(T_*\) is finite in the precise sense
\begin{equation}\label{eq:FiniteLipschitzVorticityStrainAction}
  \int_0^{T_*}
  \left(
    \norm{\nabla u(t)}_\infty
    +\norm{\omega(t)}_\infty
    +\norm{S(t)}_\infty
  \right)dt
  <\infty.
\end{equation}

\subsection{The two vortex readings}\label{sub:ClosingCriticalHeightZenoAndTheVortex}

A drawn narrowing core can refer to two different motions.  It may follow the
same material particles as they approach one another, or it may keep its visible
shape while different particles enter and leave.  These pictures are not
interchangeable, and neither exhausts possible singular geometry.  The
terminal rectangles test a precise consequence of each.

Follow the fixed labels first.
Let \(X(a,t)\) solve
\[
  \partial_tX(a,t)=u(X(a,t),t),
  \qquad
  X(a,0)=a.
\]
Two such trajectories satisfy
\[
  \frac d{dt}|X(a,t)-X(b,t)|
  \geq
  -\norm{\nabla u(t)}_\infty|X(a,t)-X(b,t)|,
\]
and the finite Lipschitz action integrates this differential inequality to
\begin{equation}\label{eq:BodyMaterialCoreLowerBound}
  |X(a,t)-X(b,t)|
  \geq
  |a-b|
  \exp\!\left(
    -\int_0^t\norm{\nabla u(s)}_\infty\,ds
  \right)
  >0
\end{equation}
Because the same row also puts \(u\) in \(L^1(I_*;L^\infty)\), every trajectory
has a limit at \(T_*\).  Letting \(t\uparrow T_*\) preserves the lower bound:
\[
  |X(a,T_*)-X(b,T_*)|
  \geq |a-b|\exp\!\left(-\int_0^{T_*}\norm{\nabla u(s)}_\infty\,ds\right)>0.
\]
Distinct labels remain spatially distinct at the terminal time.
If a drawn core collapses to one point, it cannot still contain exactly the
material core with which it began.

Even perfect centring cannot rescue a nonzero circulation from this fixed-label
picture.  The adjacent rows in \(\cR_{0,3}\), followed by
\Cref{lem:AdjacentHilbertScaleTrace}, give
\[
  u\in C([0,T_*];H^3).
\]
Together with the trace in \Cref{lem:AdjacentHilbertScaleTrace}, the embedding
\(H^3(\R^3)\hookrightarrow W^{1,\infty}(\R^3)\) gives the uniform ceiling
\begin{equation}\label{eq:BodyUniformLipschitzBound}
  L_*:=\sup_{0\leq t\leq T_*}\norm{\nabla u(t)}_\infty<\infty.
\end{equation}
Relative velocity at any chosen centre \(x_c(t)\) obeys
\[
  |u(t,x)-u(t,x_c)|\leq L_*|x-x_c|.
\]
Around a circular cross-section \(C_r(t)\), the constant centre velocity
integrates to zero.  What remains is bounded by its Lipschitz variation:
\begin{equation}\label{eq:BodyCenteredCirculationBound}
  \left|
    \oint_{C_r(t)}u(t,x)\cdot d\ell
  \right|
  \leq2\pi L_*r^2
  \longrightarrow0
  \qquad(r\downarrow0).
\end{equation}
As \(r\) shrinks, the tangential velocity relative to the centre is \(O(r)\)
and the circuit has length \(O(r)\).  Their product is \(O(r^2)\), so this
centred material core cannot carry a nonzero limiting circulation.

Now let membership change.  Particle separation no longer tracks the visible
core.  A stronger recruitment picture would keep a fixed amount of critical
content beyond successively larger frequencies.  That precise claim can be
tested in the high-frequency tail of the whole velocity field.
For \(s>1/2\) and \(K\geq1\), define
\[
  H_{>K}(t)
  :=\frac12\int_{|\xi|>K}
  |\xi|\,|\widehat u(t,\xi)|^2\,d\xi.
\]
For every \(s>1/2\), the tail satisfies
\begin{equation}\label{eq:BodyUltravioletCriticalTail}
  H_{>K}(t)
  \leq
  \frac12K^{1-2s}\norm{u(t)}_{\dot H^s}^2.
\end{equation}
Choose one integer \(m\geq s\).  The compatible assembly and adjacent trace
proved in \Cref{prop:BodyCompatibleProjectiveAssembly,lem:AdjacentHilbertScaleTrace},
followed by \(H^m\hookrightarrow\dot H^s\), make the estimate uniform in time
and force
\begin{equation}\label{eq:BodyUniformUltravioletVanishing}
  \lim_{K\to\infty}
  \sup_{0\leq t\leq T_*}H_{>K}(t)=0.
\end{equation}
Recruitment may move activity between scales, but it cannot leave a fixed amount
of critical height beyond every cutoff.  This excludes the stated persistent
tail model.  It does not prove that every geometrically narrowing core must
produce such a Fourier tail; a separate localization argument would be needed
to connect an arbitrary moving core to that frequency statement.  What is
proved is narrower and exact: fixed labels cannot collapse, their centred
circulation vanishes, and the compatible intersection cannot sustain a
cutoff-independent critical tail.

\subsection{The terminal face carries no local singular content}\label{sub:TheLocalSingularityPictureCloses}

The first rectangle has already met the global signature of a singular time:
the endpoint theorem forces the critical \(L^3\) norm to escape
\parencite{EscauriazaSereginSverak2003}, while
\eqref{eq:EndpointCriterionFromLowRectangle} gives the same velocity a finite
roof.  The complete intersection lets the same terminal face be read locally.

The mixed velocity traces and the normalized pressure jet in
\Cref{prop:PressureCompleteEndpointJet} give the terminal-uniform ceiling
\begin{equation}\label{eq:FiniteIntervalVelocityPressureSupremum}
  K_*
  :=\sup_{0\leq t\leq T_*}
  \left(
    \norm{u(t)}_{L^\infty}
    +\norm{p(t)}_{L^\infty}
  \right)
  <\infty.
\end{equation}
Fix a candidate terminal point \(z_*=(x_*,T_*)\) and place backward parabolic
cylinders at that point:
\[
  Q_\rho(z_*)
  =B_\rho(x_*)\times(T_*-\rho^2,T_*).
\]
Its spatial volume and time length contribute \(\rho^5\).  Dividing by the
critical factor \(\rho^2\) still leaves three powers of the radius:
\begin{equation}\label{eq:VanishingLocalScaleInvariantContent}
  \rho^{-2}
  \iint_{Q_\rho(z_*)}
  \left(|u|^3+|p|^{3/2}\right)\,dx\,dt
  \leq
  C\rho^3\left(K_*^3+K_*^{3/2}\right)
  \longrightarrow0
  \qquad(\rho\downarrow0).
\end{equation}
A Caffarelli--Kohn--Nirenberg singular point would retain a positive
scale-invariant amount of velocity-pressure content along a shrinking sequence
\parencite{CaffarelliKohnNirenberg1982}.  Here the terminal-uniform ceiling
makes that content vanish.

We have returned to the same vortex under the whole-terminal condition.  Its
critical height cannot complete infinitely many doublings; its stretching
action is integrable; a fixed material core cannot collapse its labels; the
stated persistent-tail recruitment model cannot survive; and every terminal
cylinder loses the velocity-pressure content a singular point would require.
These are consequences of the smooth terminal datum already used in the
same-history restart, not additional routes to an unconditional theorem.

\clearpage
\section{Further Consequences and the Euler Boundary}\label{sec:FurtherConsequencesAndEulerBoundary}

The same finite rows have two further consequences worth keeping.  One adds
the vorticity across all dyadic scales.  The other compares the Navier--Stokes
Tower with Euler and isolates the spatial payment that disappears when
viscosity is removed.

\subsection{A finite sum across infinitely many scales}\label{sub:AFiniteSumAcrossScales}

Fix \(T<T_*\).  On \([0,T]\), strict-interval parabolic regularity puts
\(u\) in \(L^2_tH^3_x\); under the whole-terminal condition and the restart
above, the same statement is available for every finite \(T\).  This one
spatial row is enough to add the vorticity from every dyadic wavelength rather
than inspect the shells one at a time.

Let \((\Delta_j)_{j\geq-1}\) be an inhomogeneous Littlewood--Paley
decomposition \parencite[Chapter~2]{BahouriCheminDanchin2011}.  At high
frequency, \(\Delta_j\omega\) sees wavelengths of order \(2^{-j}\), while its
\(L^\infty\) norm records the largest rotation in that shell.  Define
\[
  \norm{\omega}_{B^0_{\infty,1}}
  :=\sum_{j\geq-1}\norm{\Delta_j\omega}_\infty.
\]
Bernstein's inequality costs three halves of a derivative in passing from
\(L^2\) to \(L^\infty\):
\[
  \norm{\Delta_j\omega}_\infty
  \leq C2^{3j/2}\norm{\Delta_j\omega}_2.
\]
The \(H^2\) weight of \(\omega\) supplies two derivatives.  Splitting
\(2^{3j/2}=2^{-j/2}2^{2j}\) exposes the spare half derivative, and
Cauchy--Schwarz sums it:
\begin{align*}
  \sum_{j\geq0}\norm{\Delta_j\omega}_\infty
  &\leq
  C\left(\sum_{j\geq0}2^{-j}\right)^{1/2}
  \left(\sum_{j\geq0}2^{4j}\norm{\Delta_j\omega}_2^2\right)^{1/2}
  \\
  &\leq C\norm{\omega}_{H^2}
  \leq C\norm{u}_{H^3}.
\end{align*}
The low shell is controlled directly by \(\norm{\omega}_2\).  Thus one finite
Sobolev coordinate pays the full \(\ell^1\) shell sum:
\begin{equation}\label{eq:VorticityBesovFromHThree}
  \norm{\omega}_{B^0_{\infty,1}}
  \leq C\norm{u}_{H^3}.
\end{equation}
Cauchy--Schwarz in time now turns the \(L^2_tH^3_x\) row into
\begin{equation}\label{eq:FiniteDyadicVorticityAction}
  \sum_{j\geq-1}
  \int_0^{T}\norm{\Delta_j\omega(t)}_\infty\,dt
  =\int_0^{T}\norm{\omega(t)}_{B^0_{\infty,1}}\,dt
  <\infty.
\end{equation}
Infinitely many shells may still become active.  What cannot happen on the
fixed classical interval is a nonsummable succession of their time-integrated
peaks.  The conclusion uses one finite Sobolev row; it asserts neither an
analytic radius nor a weighted estimate over all derivative orders, and its
extension to every finite time follows only after the conditional restart has
made the history global.

\Needspace{8\baselineskip}
\subsection{Why the Euler Tower is different}\label{sub:WhyTheEulerTowerIsDifferent}

Remove viscosity while keeping the two histories separate.  Both equations can
still be differentiated in time, so both possess a mixed jet.  Write
\[
  V_n^{\mathrm{NS}}:=\partial_t^nu^{\mathrm{NS}},
  \qquad
  Q_n^{\mathrm{NS}}:=\partial_t^np^{\mathrm{NS}},
\]
and define \(V_n^{\mathrm E},Q_n^{\mathrm E}\) from an Euler solution with its
own history.  The two recurrences differ in one place:
\begin{align}
  V_{n+1}^{\mathrm{NS}}
  &+
  \sum_{a=0}^{n}\binom{n}{a}
  (V_a^{\mathrm{NS}}\cdot\nabla)V_{n-a}^{\mathrm{NS}}
  +\nabla Q_n^{\mathrm{NS}}
  =\nu\Delta V_n^{\mathrm{NS}},
  \label{eq:NSTemporalRecurrenceComparison}
  \\
  V_{n+1}^{\mathrm E}
  &+
  \sum_{a=0}^{n}\binom{n}{a}
  (V_a^{\mathrm E}\cdot\nabla)V_{n-a}^{\mathrm E}
  +\nabla Q_n^{\mathrm E}
  =0.
  \label{eq:EulerTemporalRecurrenceComparison}
\end{align}
Transport and the nonlocal pressure response remain in both equations.
Only the Navier--Stokes row contains \(\nu\Delta V_n^{\mathrm{NS}}\), which
connects a temporal response to two more spatial derivatives of that same
response.

The energy rows expose the same difference.  Set
\[
  E_s^{\mathrm{NS}}:=\frac12\norm{\Lambda^su^{\mathrm{NS}}}_2^2,
  \qquad
  E_s^{\mathrm E}:=\frac12\norm{\Lambda^su^{\mathrm E}}_2^2,
\]
with the corresponding nonlinear transfer terms.  Then
\begin{equation}\label{eq:NSEulerEnergyComparison}
  (E_s^{\mathrm{NS}})'
  =\mathcal P_s^{\mathrm{NS}}-2\nu E_{s+1}^{\mathrm{NS}},
  \qquad
  (E_s^{\mathrm E})'
  =\mathcal P_s^{\mathrm E}.
\end{equation}
Repeated time differentiation of the Navier--Stokes row, followed by division
by \((-2\nu)^n\), isolates \(E_{n+1/2}^{\mathrm{NS}}\) in
\eqref{eq:BodyCompletedCriticalRow}.  The Euler row still mixes time and space
nonlinearly, but no signed coercive term singles out the next spatial height.

If an Euler history were independently known to lie in the analogous
all-orders intersection, then every fixed Sobolev norm would be controlled and
in particular it would satisfy the Beale--Kato--Majda condition
\(\int_0^T\|\omega^{\mathrm E}(t)\|_\infty\,dt<\infty\)
\parencite{BealeKatoMajda1984}.  The present theorem supplies no such Euler
intersection, and none of its constants is asserted to remain bounded as
\(\nu\downarrow0\).

Euler keeps the jet, transport, pressure, and their nonlinear coupling.  It
loses the viscous rung that places an adjacent higher spatial energy in the
Navier--Stokes row.  At fixed \(\nu>0\), that rung supplies the strict-interval
parabolic gain and exposes the terminal rectangles.  If the whole-terminal
condition holds, those rectangles reach one datum and the same Navier--Stokes
history restarts globally with bounded kinetic energy.  What this manuscript
does not derive for arbitrary large \(u_0\) is the first cutoff-uniform payment
\(D_*<\infty\); without it, neither the lowest terminal rectangle nor the
global conclusion has been proved.
